\section{Synchronous machine --- \texttt{genqec}}
\label{sec:genqec}

\texttt{genqec} is the round-rotor synchronous machine model that PSLF introduced
as the successor to \texttt{genrou} and \texttt{gentpj}. It retains the two-axis,
two-rotor-circuit-per-axis structure of \texttt{genrou} (transient and subtransient
time constants and reactances on both axes) but represents magnetic saturation on the
equivalent circuit rather than by scaling the exciting flux alone, which gives a
better match to measured field current at load. The saturation function is quadratic
through the points $S(1.0)$ and $S(1.2)$; \texttt{kw} weights the saturation between
the axes and \texttt{satflg} selects the implementation, and both are left at their
defaults of zero here. The compounding terms \texttt{rcomp}/\texttt{xcomp} are zero,
so the voltage regulator sees the unloaded terminal voltage.

The two machine families differ in most electrical constants, Table~\ref{tab:genqec}.
The Baker Hughes units have a considerably longer $T'_{d0}$ (10\,s against 5.75\,s),
higher synchronous and transient reactances and slightly larger inertia. For each
family the ordering $L_d > L'_d > L''_d \geq L_l$ and $L_q > L'_q \geq L''_q$ holds and
$T'_{d0} > T''_{d0}$, $T'_{q0} > T''_{q0}$, so the data are internally consistent.
The inertia is given on the machine base; on the 100\,MV$\cdot$A system base the
Turboden units contribute $H_{sys} = 1.49 \cdot 76.471 / 100 = 1.139$\,s and the Baker
Hughes units $H_{sys} = 1.50 \cdot 71.13 / 100 = 1.067$\,s each.

\begin{table}[!ht]
\centering
\begin{tabularx}{\linewidth}{@{}l X l r r@{}}
\toprule
\textbf{Parameter} & \textbf{Description} & \multicolumn{1}{c}{\textbf{Unit}} & \multicolumn{1}{c}{\textbf{Turboden}} & \multicolumn{1}{c}{\textbf{Baker Hughes}} \\
\midrule
\texttt{mva} & Machine MVA base & MV$\cdot$A & 76.471 & 71.13 \\
\texttt{tpdo} & d-axis transient open-circuit time constant $T'_{d0}$ & s & 5.751 & 10 \\
\texttt{tppdo} & d-axis subtransient open-circuit time constant $T''_{d0}$ & s & 0.044 & 0.05 \\
\texttt{tpqo} & q-axis transient open-circuit time constant $T'_{q0}$ & s & 0.865 & 0.865 \\
\texttt{tppqo} & q-axis subtransient open-circuit time constant $T''_{q0}$ & s & 0.072 & 0.13 \\
\texttt{h} & Inertia constant $H$ & MW$\cdot$s/MV$\cdot$A & 1.49 & 1.5 \\
\texttt{d} & Damping factor $D$ & pu & 3 & 3 \\
\texttt{ld} & d-axis synchronous reactance $L_d$ & pu & 1.436 & 1.49 \\
\texttt{lq} & q-axis synchronous reactance $L_q$ & pu & 0.61 & 0.78 \\
\texttt{lpd} & d-axis transient reactance $L'_d$ & pu & 0.189 & 0.21 \\
\texttt{lpq} & q-axis transient reactance $L'_q$ & pu & 0.61 & 0.78 \\
\texttt{lppd} & d-axis subtransient reactance $L''_d$ & pu & 0.121 & 0.14 \\
\texttt{lppq} & q-axis subtransient reactance $L''_q$ & pu & 0.181 & 0.18 \\
\texttt{ll} & Stator leakage reactance $L_l$ & pu & 0.0991 & 0.11 \\
\texttt{s1} & Saturation factor at 1.0 pu flux $S(1.0)$ & -- & 0.1 & 0.1 \\
\texttt{s12} & Saturation factor at 1.2 pu flux $S(1.2)$ & -- & 0.522 & 0.522 \\
\texttt{ra} & Stator resistance $R_a$ & pu & 0.002 & 0.0045 \\
\texttt{rcomp} & Compounding resistance (voltage-regulator compensation) & pu & 0 & 0 \\
\texttt{xcomp} & Compounding reactance (voltage-regulator compensation) & pu & 0 & 0 \\
\texttt{kw} & Saturation weighting factor $K_w$ & -- & 0 & 0 \\
\texttt{satflg} & Saturation-function selector & -- & 0 & 0 \\
\bottomrule
\end{tabularx}

\caption{\texttt{genqec} parameters (pu on machine base unless noted).}
\label{tab:genqec}
\end{table}
